% Physical unit 009; accepted Session 02.
\GordanUnitStart{009}{printed page 4}
\GordanTopFolio{4}
\begin{center}
{\large CHAPTER II.}\\[0.8em]
\S\ 1.
\end{center}
A curve described by the point $\Pi$ on the spheroid, while no force acts on it except a normal force, is called a geodetic line. We find its equation by eliminating from the equations of motion of the moving point the quantities $t$ and $\lambda$, of which the latter represents the normal force and the former a proportional time.

\begin{center}\S\ 2.\end{center}
Taking, therefore, the equation of the spheroid as \gq{$f=0$}, the equations of motion in rectangular coordinates are
\[
 {d_2x\over dt^2}=\lambda{\partial f\over\partial x},\qquad
 {d_2y\over dt^2}=\lambda{\partial f\over\partial y},\qquad
 {d_2z\over dt^2}=\lambda{\partial f\over\partial z}.
\]
If these are successively multiplied by the differentials
\[
 dx\qquad dy\qquad dz,
\]
and added, there results
\[
 {dx\over dt}{d_2x\over dt^2}+{dy\over dt}{d_2y\over dt^2}+{dz\over dt}{d_2z\over dt^2}=0;
\]
or, on integrating,
\[
 {dx\over dt}{dx\over dt}+{dy\over dt}{dy\over dt}+{dz\over dt}{dz\over dt}={1\over C^2};
\]
that is,
\[
 dt=C\,ds.
\]
Thus we see the familiar fact that the velocity of the point $\Pi$ is constant.

\begin{center}\S\ 3.\end{center}
Introducing this value into the equations of motion, we have
\[
 {d_2x\over ds^2}=\lambda C^2{\partial f\over\partial x},\qquad
 {d_2y\over ds^2}=\lambda C^2{\partial f\over\partial y},\qquad
 {d_2z\over ds^2}=\lambda C^2{\partial f\over\partial z}.
\]
These express nothing else than this: the variation of the extension, while the relation $f=0$ is maintained in forming it, itself vanishes.

Therefore the equation of the geodetic line is
